\documentclass[11pt]{article}
\usepackage[a4paper,margin=25mm]{geometry}
\usepackage{amsmath,amssymb,amsthm,booktabs,hyperref}
\newtheorem{theorem}{Theorem}
\newtheorem{lemma}[theorem]{Lemma}
\newtheorem{corollary}[theorem]{Corollary}
\newcommand{\F}{\mathbb F}
\newcommand{\GL}{\operatorname{GL}}
\newcommand{\SL}{\operatorname{SL}}
\newcommand{\SU}{\operatorname{SU}}
\newcommand{\PSL}{\operatorname{PSL}}
\newcommand{\PGL}{\operatorname{PGL}}
\newcommand{\Aut}{\operatorname{Aut}}
\newcommand{\gen}[1]{\langle #1\rangle}
\title{Short positive generation certificates in the acquired $G_2(11)$ model}
\author{C5 (\texttt{/root/c5\_terminal\_audit})}
\date{20 October 2026}
\begin{document}
\maketitle

\begin{abstract}
The complete maximal-subgroup classification implies that a subgroup of
$G_2(11)$ containing elements of orders $19$ and $37$ is the whole group.
No irreducibility hypothesis is needed.  Moreover, a proper subgroup
containing an element of order $37$ preserves a nonisotropic line in the
natural seven-dimensional Cayley module.  An order-$37$ matrix itself
has exactly one rational invariant line, so this gives a second positive
certificate using one word, one fixed vector, and one nonzero minor.
We also give the full sign-eigenspace formulation for involutory marks.
We record the exceptional-subgroup checks, exact witness duties, and the
limited relation to the existing generation and intersection tests.
No finite-field, group, matrix, Reader, control or census computation was
performed for this note; no candidate is decided by it.
\end{abstract}

\section{The complete classification input}

Put $G=G_2(11)$, the finite split Chevalley group, without a covering or
semilinear extension.  The classification used here is Kleidman's
classification for odd $q$~\cite{K}.  Its complete list for characteristic
at least five is explicitly reproduced before the subsequent character
degree elimination in Nguyen~\cite[proof of Theorem 4.4, pp.~10--11]{N}.
It is important to use that full list, rather than the shorter conclusion
of Theorem 4.4, whose cross-characteristic irreducibility hypotheses are
irrelevant here.  Wilson's retained book independently records the full
list in Table~4.1~\cite{W}.

At $q=11$ every maximal subgroup has one of the following isomorphism
types.  Rows refer to types, not to an assertion that there is just one
conjugacy class of each type.  Their orders and prime divisors suffice
for the argument.

\begin{center}
\small
\begin{tabular}{lll}
\toprule
Type & Order & Prime divisors\\
\midrule
Either maximal parabolic & $11^6(11-1)^2(11+1)$ & $2,3,5,11$\\
$(\SL_2(11)\circ\SL_2(11)).2$
 & $11^2(11^2-1)^2$ & $2,3,5,11$\\
$\SL_3(11):2$ & $2\cdot11^3(11^3-1)(11^2-1)$
 & $2,3,5,7,11,19$\\
$\SU_3(11):2$ & $2\cdot11^3(11^3+1)(11^2-1)$
 & $2,3,5,11,37$\\
$2^3.\PSL_3(2)$ & $2^6\cdot3\cdot7$ & $2,3,7$\\
$\PGL_2(11)$ & $11(11^2-1)$ & $2,3,5,11$\\
$G_2(2)\cong\operatorname{PSU}_3(3):2$
 & $2^6\cdot3^3\cdot7$ & $2,3,7$\\
$J_1$ & $2^3\cdot3\cdot5\cdot7\cdot11\cdot19$
 & $2,3,5,7,11,19$\\
\bottomrule
\end{tabular}
\end{center}

Here $\circ$ identifies the central involutions of the two $\SL_2$
factors.  Thus taking that central product and adjoining the outer factor
of order two gives exactly the displayed involution-centralizer order.
The two parabolic radicals each have order $11^5$; their common Levi order
is $|\GL_2(11)|=11(11-1)^2(11+1)$.  The linear and unitary entries are
the full $\SL_3$ and $\SU_3$ entries, including their centres.  In
particular the unitary centre is not discarded in the divisibility check.
Wilson's Section~5.9.1 supplies the stated $J_1$ order and its embedding
in $G_2(11)$~\cite{W}; the exceptional $J_1$ case really contains the
prime $19$, but not $37$.

For completeness, all field and exceptional restrictions in the general
classification have been accounted for.  There is no proper subfield
subgroup because $11$ is prime.  Characteristic-three cases, including
the Ree subgroups and additional characteristic-three classes, do not
apply.  The characteristic-two $J_2$ cases do not apply.  The remaining
general exceptional possibilities $\PSL_2(13)$ and $\PSL_2(8)$ do not
occur in this prime field: the prime-field restrictions are respectively
$p\equiv1,3,4,9,10,12\pmod {13}$ and $p\equiv1,8\pmod9$
in Burness--Thomas~\cite[Proposition 7.4(a,b)]{BT}, and $p=11$
satisfies neither.  Even retaining the larger envelopes
$\PGL_2(13)$ and $\PSL_2(8):3$ would introduce no prime $19$ or $37$,
so the two conclusions below do not depend on excluding those cases.
No exceptional almost-simple case has been silently removed by a
natural-module irreducibility assumption.

The original publisher full text of~\cite{K} was not retrieved during
this assessment.  The complete statement was checked in the retained
Wilson book and the displayed list in the primary research paper~\cite{N};
the exceptional field restrictions were checked in the authors' manuscript
of~\cite{BT}.  The classification remains an attributed theorem input,
not a classification proved anew by this note.

\section{Two prime witnesses}

\begin{theorem}\label{thm:two}
If $H\leq G_2(11)$ contains an element of order $19$ and an element of
order $37$, then $H=G_2(11)$.  No assumption on a marking, intersection
property, or irreducibility of $H$ is required.
\end{theorem}
\begin{proof}
If $H<G$, finiteness puts $H$ inside a maximal subgroup $M$.
Lagrange's theorem would make both $19$ and $37$ divide $|M|$.
No row of the complete table has both primes.  This is a contradiction.
\end{proof}

The relevant elementary factorizations are
\[
 11^3-1=(11-1)(11^2+11+1)=2\cdot5\cdot7\cdot19,
 \qquad
 11^3+1=(11+1)(11^2-11+1)=2^2\cdot3^2\cdot37.
\]
The primitive-divisor language motivates the choice, but the proof needs
only the explicit classification and divisibility table.  It does not
assume a general primitive-prime generation theorem with unchecked
exceptions.  Neither one of the two primes suffices: the proper
$\SL_3(11):2$ and $\SU_3(11):2$ entries have the respective primes in
their orders, and Cauchy's theorem supplies such elements.

\begin{corollary}[Prime-order word certificate]\label{cor:words}
Let $h_1,\ldots,h_r$ be actual members of an acquired matrix group
$M\cong G_2(11)$, and $H=\gen{h_1,\ldots,h_r}$.  Suppose explicitly
specified words evaluate to $X,Y$ and satisfy
\[
 X\neq I,\quad X^{19}=I,\qquad
 Y\neq I,\quad Y^{37}=I.
\]
Then $H=M$.
\end{corollary}
\begin{proof}
Each word belongs to $H$.  Since $19$ and $37$ are prime, the displayed
conditions prove their exact respective orders.  Apply
Theorem~\ref{thm:two} through the acquired isomorphism.
\end{proof}

Thus no exact order algorithm or complete subgroup closure is needed to
verify a supplied positive certificate.  If instead an actual word $w$
has certified exact order $n$ divisible by a prime $p$, then
$w^{n/p}$ has order $p$, so the older order-divisibility certificate is
also sound.  An unproved multiple of the order cannot silently replace
the exact order in that inference.  Directly checking the prime-power
equalities in Corollary~\ref{cor:words} avoids this issue.

\begin{lemma}[Adjacent products lose no positive certificate]\label{lem:even}
If $h_1,\ldots,h_r$ are involutions, let
$H^+=\gen{h_1h_2,h_2h_3,\ldots,h_{r-1}h_r}$.
Then $H^+\trianglelefteq H$ and $[H:H^+]\leq2$.
If $H\cong G_2(11)$, then $H^+=H$.
\end{lemma}
\begin{proof}
Every pair product $h_i h_j$ is a product of adjacent products or their
inverses, by telescoping.  Conjugation by any $h_k$ takes $h_i h_j$ to
$(h_kh_i)(h_jh_k)$, again a product of pair products.  In the quotient
all $h_i$ have the same image, of order at most two.  The final claim
uses that $G_2(11)$ is nonabelian simple and has no quotient of order two.
\end{proof}

The lemma explains why words in adjacent products are a sound search
domain.  A bounded list of $160$ letters, or any other capped search,
need not find the witnesses even when $H=G$.  A failed search proves
neither properness nor nongeneration.  No word-length or runtime guarantee
is asserted here.

\section{A second sufficient certificate: order $37$ and no invariant line}

The following elementary maximal-overgroup observation explains the
certificate architecture.  It is a general structure-stabilizer argument,
not a new classification theorem.

\begin{lemma}[Unique fixed structure]\label{lem:structure}
Let a finite group $K$ act on a set $\mathcal X$.  Let $w\in K$ fix
exactly one point $x\in\mathcal X$, and suppose every maximal proper
subgroup containing $w$ fixes some point of $\mathcal X$.  For any
$\gen{w}\leq H\leq K$, properness implies $H\leq K_x$.
Consequently one $H$-element moving $x$ proves $H=K$.
If also $K_x<K$, then
\[
 H<K\quad\Longleftrightarrow\quad H\leq K_x.
\]
For a specified generating set of $H$, these alternatives are certified
by every generator fixing $x$, or by one generator moving $x$.
\end{lemma}
\begin{proof}
A proper $H$ lies in a maximal proper $L$.  If $L$ fixes $y$, then its
element $w$ fixes $y$, and uniqueness forces $y=x$.  Hence $H\leq K_x$.
The reverse implication in the final assertion uses the additional
hypothesis that $K_x$ is proper.  That hypothesis cannot be omitted
from the biconditional in an arbitrary group action.
\end{proof}

Let $C$ be a split Cayley algebra over $\F_{11}$ and let $V=C_0$ be its
seven-dimensional trace-zero space.  Work in the actual natural action
$G=\Aut_{\F_{11}}(C)$ on $V$.

\begin{theorem}\label{thm:line}
Every proper subgroup $H<G$ containing an element of order $37$
preserves a nonisotropic one-dimensional $\F_{11}$-subspace of $V$.
Consequently, an order-$37$ word and a certificate that $H$ preserves no
common $\F_{11}$-line prove $H=G$.
\end{theorem}
\begin{proof}
Choose a maximal overgroup of a proper $H$.  The table leaves only
$\SU_3(11):2$, since it is the only maximal type whose order is divisible
by $37$.  In the natural Cayley action this subgroup is the stabilizer
of a nonisotropic line of $C_0$.  Wilson proves that geometric
description in Section~4.2.6 of the retained edition~\cite{W}.
Every conjugate has the same property.  Thus $H$ preserves that line.
The contrapositive proves the certificate assertion.
\end{proof}

This result needs the natural Cayley-module identification; an arbitrary
faithful representation of the abstract group is not enough to infer
the indicated line stabilizer.  No full seven-dimensional
irreducibility test is required.  It is sufficient, and a little stronger
than necessary, to exclude all common rational lines, including isotropic
ones.

\subsection*{The smallest certificate uses the order-$37$ word itself}

\begin{lemma}\label{lem:unique}
Let $Y\in\GL_7(\F_{11})$ have order $37$.  Then its fixed space has
dimension one and is its unique invariant one-dimensional
$\F_{11}$-subspace.
\end{lemma}
\begin{proof}
The order of $11$ in $(\mathbb Z/37\mathbb Z)^\times$ is six:
$11^2\equiv10$, $11^3\equiv-1$, and none of $11,11^2,11^3$ is one
modulo $37$.  Frobenius therefore has orbits of size six on all nonidentity
$37$th roots of unity.  Consequently every nontrivial irreducible factor
of $X^{37}-1$ over $\F_{11}$ has degree six; the remaining factor is
$X-1$.  The polynomial is separable since the characteristic is not $37$.

Since $Y^{37}=I$, the module is a direct sum of the corresponding simple
constituents.  Since $Y\neq I$, at least one constituent is nontrivial
and has dimension six.  Total dimension seven then leaves exactly one
trivial constituent and cannot contain another nontrivial constituent.
Thus $\dim\ker(Y-I)=1$.  Finally, on a rational invariant line $Y$ acts
by a scalar $\lambda\in\F_{11}^{\times}$ with
$\lambda^{37}=\lambda^{10}=1$.  Coprimality gives $\lambda=1$, so that
line must be the fixed line just found.
\end{proof}

\begin{corollary}[One word, one vector, one moved line]\label{cor:moved}
In the natural acquired $G_2(11)$ module, let
$H=\gen{h_1,\ldots,h_r}\leq G$.  Suppose a recorded $H$-word evaluates
to $Y\neq I$ with $Y^{37}=I$.  Suppose a nonzero row $v$ satisfies
$vY=v$, and for at least one actual generator $h_j$ the rows $v$ and
$vh_j$ are linearly independent.  Then $H=G$.
\end{corollary}
\begin{proof}
The lemma makes $\gen{v}$ the only rational line invariant under $Y$.
Every $H$-invariant line must be such a line, because $Y\in H$.
The displayed generator moves this line, so no $H$-invariant rational
line exists.  Theorem~\ref{thm:line} now applies.
\end{proof}

No involution assumption on the $h_j$ is needed for this corollary.
Linear independence can be witnessed by two coordinate positions $s,t$
for which
\[
 v_s(vh_j)_t-v_t(vh_j)_s\neq0.
\]
Together with the prime-order word, the seven fixed-vector equalities,
and $v\neq0$, this is a short fully exact sufficient certificate.
It requires no second prime word, complete subgroup closure, full
irreducibility computation, enumeration of rational lines, or exhaustive
sign systems.  The fixed space can be found by ordinary linear solving
in a future acquired implementation.  Its uniqueness here is a proved
consequence of the order and dimension, not an inference from one
found vector.  In a different dimension or over a different field this
argument must be re-established; it is not a general fixed-line heuristic.
If no order-$37$ word or moved-line witness is obtained, this certificate
has not by itself proved generation.  There is also a complete positive
properness certificate after all generator-line checks are completed,
as follows.

\begin{corollary}[Exact dichotomy once the word is supplied]\label{cor:dichotomy}
Under the hypotheses of Corollary~\ref{cor:moved} concerning $Y$ and $v$,
the fixed line $\gen{v}$ is nonisotropic and its full stabilizer in $G$
is a proper subgroup of type $\SU_3(11):2$.  Thus
\[
 H<G\quad\Longleftrightarrow\quad
 \text{every }h_i\text{ preserves }\gen{v}.
\]
\end{corollary}
\begin{proof}
The cyclic subgroup $\gen{Y}$ is proper in the noncyclic $G$.  Any
maximal overgroup has the unitary type in the table and preserves a
nonisotropic line.  Lemma~\ref{lem:unique} makes that line $\gen{v}$.
Wilson's geometric description identifies its full stabilizer as the
proper unitary normalizer.  Apply Lemma~\ref{lem:structure}.
\end{proof}

For involutory generators, a preserved line satisfies $vh_i=v$ or
$vh_i=-v$, and one may record that sign for each generator.  For arbitrary
generators the proportionality scalar is simply a nonzero field element.
An unfinished search for $Y$, an unverified proposed fixed vector, or
unfinished generator checks remain inconclusive.  The dichotomy does
not turn earlier partial or sufficient-only line-finder returns into
completed decisions.

\subsection*{An independent full line-absence certificate}

For involutory generators the latter certificate is small and exact.
Use the acquired right-row convention $x\mapsto xh_i$.  For every
$\epsilon=(\epsilon_1,\ldots,\epsilon_r)\in\{1,-1\}^r$ form
\[
 A_\epsilon=
 [\,h_1-\epsilon_1 I\mid h_2-\epsilon_2 I\mid\cdots
       \mid h_r-\epsilon_r I\,],
 \qquad A_\epsilon\in\operatorname{Mat}_{7\times7r}(\F_{11}).
\]
There is no common rational line if and only if every $A_\epsilon$ has
rank seven.  Indeed, a common line has a nonzero row generator $x$;
on it every involution acts by a scalar whose square is one, hence by
one of the two indicated signs.  Conversely a nonzero common row
nullvector supplies such a line.  A nonzero $7\times7$ minor for each
sign pattern certifies all these ranks.  There are $8$, $16$ or $32$
systems for respectively three, four or five generators.  This is a
written certificate specification, not an executed rank test.

The normalizer factor of order two must be retained in this argument.
Although the unitary subgroup fixes an appropriate vector, its full
$\SU_3(11):2$ normalizer may send that vector to its negative.  Checking
only common fixed vectors (all signs positive) is therefore insufficient.
Checking every indicated sign system handles the full line stabilizer.
Likewise the proof does not depend on identifying all conjugacy classes
of that type with each other: every surviving maximal embedding has the
specified natural line-stabilizer property.

\subsection*{Reusing the acquired ordered endpoint decomposition}

There is a further exact linear reduction when an actual commuting
endpoint pair $(a,e)$ has its acquired simultaneous pure eigenspaces
\[
 V_{\alpha,\eta}=\{x:xa=\alpha x,\ xe=\eta x\},
 \qquad (\alpha,\eta)\in\{1,-1\}^2,
\]
of dimensions $1,2,2,2$ in the stated order.  Suppose $B_{\alpha,\eta}$
is a full row-basis matrix of that actual space, of row rank
$d_{\alpha,\eta}$.  For a rank-four candidate $(a,b,c,e)$, a common
invariant line exists if and only if, for at least one of the sixteen
choices $(\alpha,\eta,\sigma,\tau)$, the matrix
\[
 [\,B_{\alpha,\eta}(b-\sigma I)
        \mid B_{\alpha,\eta}(c-\tau I)\,]
 \quad\text{of size }d_{\alpha,\eta}\times14
\]
has row rank smaller than $d_{\alpha,\eta}$.  This follows by writing
$x=zB_{\alpha,\eta}$; full row rank makes $x\neq0$ equivalent to
$z\neq0$.  Certifying full row rank for every case therefore proves
absence, with minors of order at most two after the endpoint basis
joins.  The same argument restricts any other selected generators.
Neither a full quaternion frame nor the $14763$-element catalogue is
needed just to check these restricted linear systems.

An existing endpoint decomposition is reusable only for the exact
ordered pair, its transported actual seven-space, and its whole basis
and dimension certificates.  A decomposition of a different endpoint
pair or an unproved $1/2/2/2$ dimension label supplies no such basis.
One can always use the unreduced $7\times7r$ systems instead.

Known word relations may reduce the sign patterns further, with an
explicit proof for every discarded pattern.  For example, if
$|h_i h_j|$ is certified odd, a common line would satisfy
$(\epsilon_i\epsilon_j)^{|h_i h_j|}=1$, whence
$\epsilon_i=\epsilon_j$.  The complete $2^r$ route needs no such
optimization.  Discarding a sign merely because the generators are
adjacent, or because a prospective group is hoped to be perfect, is not
valid before generation is proved.

The earlier candidate3610 finder/checker diagnostics are not such a
certificate.  They concern specified end triples and sufficient
null-line witnesses (including the stated dual construction); their
retained \texttt{NO\_PROPOSAL\_UNPAID} outcomes cannot be promoted to
nonexistence of every invariant rational line of a new whole marking.

\section{Existing terminal predicates and the actual reusable delta}

This prime pair is already present in the acquired scientific Sources.
The frozen \texttt{matrix\_end\_certificate.g} function
\texttt{F004CertifiesFullEndMatrix} takes exactly three known-member
marks, lifts them through the acquired frame, searches words in the
two adjacent products, and returns true only after observed exact matrix
orders have both required prime divisors.  At $q=11$ its primes are
$[19,37]$.  The frozen \texttt{interval\_certificate\_v1.g} uses the
same pair for arbitrary contiguous intervals, currently with permutation
words; length three dispatches to the end collaborator.  The theorem
above grounds that existing attributed maximal-subgroup premise.  It is
not a newly acquired shortcut or evidence that a pending case succeeds.

There is a narrow possible future integration: supplied prime-order
word matrices can certify a full four- or five-mark group through the
same faithful bridge, without closing its large permutation action.
This extends the positive-certificate domain beyond the existing
three-mark matrix collaborator.  It would require a reviewed binding
and actual witnesses; no Source function, caller, control or gate has
been changed here.  The optional order-$37$/line-rank certificate is
likewise new only as a specified sufficient leg, not an implemented or
acquired predicate.  All unsuccessful certificate searches must retain
the existing exact fallback or remain unpaid.

For a rank-four marking $(a,b,c,e)$ put
\[
 L=\gen{a,b,c},\quad R=\gen{b,c,e},\quad D=\gen{b,c}.
\]
The intersection property forces $L,R<G$.  Hence a positive full-group
certificate on either end correctly rejects that marking.  A positive
certificate on all four marks proves only full generation.  It does
not prove the string commutations, properness of both ends, either
end's rank-three intersection property, or $L\cap R=D$.
These are the distinct predicates retained by
\texttt{candidate\_inventory\_no\_exact\_product.g}.

The existing rank-three test can remain genuinely small in formulation.
For distinct involutions $(s,t,u)$ satisfying $su=us$, put
$m=|st|$, $E=\gen{t,u}$ and $r=st$.  Then
\[
 \gen{s,t}\cap E=\gen{t}
 \quad\Longleftrightarrow\quad
 r^{m/p}\notin E\ \text{for every prime }p\mid m.
\]
For if the intersection exceeds $\gen{t}$, multiplying an outside
reflection by $t$ if necessary supplies a nontrivial rotation.  A
nontrivial subgroup of the cyclic rotation group contains an element
of prime order, whose subgroup is generated by one of the displayed
powers.  The converse is immediate, since a nonidentity rotation cannot
belong to $\gen{t}$.  This is exactly the specialization in the acquired
\texttt{dihedral\_intersection.g}; its reversal chooses the smaller
adjacent product order.  Exact product orders and the membership tests
remain duties.  The formulation contains no order calculation on the
whole rank-three end, but makes no measured cost or fit promise.

For rank five, use $(a,b,c,d,e)$, with proper rank-four ends
$P=\gen{a,b,c,d}$ and $Q=\gen{b,c,d,e}$ and middle
$M_0=\gen{b,c,d}$.  A positive certificate on $P$ or $Q$ is an immediate
properness obstruction.  Passing generation on all five marks still
does not establish that the two marked ends are string C-groups or
$P\cap Q=M_0$.  The outer triples $\gen{a,b,c}\leq C_G(e)$ and
$\gen{c,d,e}\leq C_G(a)$ have the still stronger pre-existing restriction
that their orders divide $11^2(11^2-1)^2$.  A witnessed element there
whose order has a prime factor $7$, $19$ or $37$ already contradicts the
required commutations in the identified group.  This is the retained
centralizer obstruction, not a universal rank-five nonexistence proof.

Finally, the complete $14763$-involution centralizer catalogue and the
$200/288$ endpoint domains of the separate inventory theorem replace
inventory construction.  They do not supply these terminal generation
or intersection results.  Both catalogue completeness and all original
prefix/index/mark transport obligations remain separate.

\section{Certificate duties and proof boundary}

A future application must authenticate all of the following.
\begin{enumerate}
\item The actual matrices are in the declared $M\cong G_2(11)$, and the
same group is the intended saved group under the acquired faithful frame.
For a catalogue-created matrix, algebra-automorphism membership can be
transported only after the full $M=\Aut(C)$ bridge is retained.
Faithfulness or preservation of one bilinear form alone does not identify
the ambient group.
\item The exact original prefix, positive old index, ordered marks,
export and native closure, and the common right-row convention are joined.
Words use the indicated actual subgroup generators, with recorded
inverses/product order.  An abstract element of the ambient group is not
automatically a word of the subgroup being certified.
\item Retain the complete words (or explicit finite straight-line
programs), evaluated matrices, and all nonidentity/power checks.  In the
line version also retain every sign system and nonzero rank minor.
Every claimed equality is an exact equality over $\F_{11}$, without
floating point or an unproved projective scalar quotient.
\item Root's genuine execution, full current Source/data/runtime/alias
authentication, limits, output inventory, closure, distinct review and
adoption remain owed.  This document neither invokes a verifier nor
declares a hypothetical witness acquired.
\end{enumerate}

The acquired q11 normalization gate
\texttt{208431dc9dd0ec31b7688719e12def55682dd35a085f088ba2cd32d34fdde940},
forms gate
\texttt{0ab1b853480874ef685a807d5c72d5f98fba22894e7b5727e7cb6815d0cc1d36}
and faithful-frame gate
\texttt{a8746058d9d1482ac9ed92157f31b765c0998d195efc312d734a8533ddc86498}
have distinct duties.  The written full-group bridge is
\texttt{acquired\_cayley\_full\_group\_bridge.rst}, digest
\texttt{a754edb570ef9d30816c3897deeae2aec7958fb7cfb5fa756536ca8e0ab30801}.
These pins are retained locators, not newly replayed acquisitions.

This is a classical classification corollary and a small certificate
design for the existing intake.  There is no claimed historical novelty,
new implementation, actual word/rank witness, closed candidate, complete
finite domain, census count, whole-original recovery, or broader
conjecture result.  The existing consumed Sources remain unchanged.

\begin{thebibliography}{9}
\bibitem{K} P.~B.~Kleidman,
The maximal subgroups of the Chevalley groups $G_2(q)$ with $q$ odd,
the Ree groups ${}^2G_2(q)$, and their automorphism groups,
\emph{J. Algebra} \textbf{117} (1988), 30--71.
\href{https://doi.org/10.1016/0021-8693(88)90239-6}{doi:10.1016/0021-8693(88)90239-6}.

\bibitem{N} H.~N.~Nguyen,
Irreducible restrictions of Brauer characters of the Chevalley group
$G_2(q)$ to its proper subgroups, primary manuscript,
\href{https://arxiv.org/pdf/0910.4994}{arXiv:0910.4994v1},
proof of Theorem~4.4, printed pp.~10--11.  The complete list occurs
before its degree-based exclusions; those exclusions are not used here.

\bibitem{W} R.~A.~Wilson, \emph{The Finite Simple Groups},
retained preliminary text \texttt{wilson\_early\_book.txt},
633091 bytes, SHA256
\texttt{d0af6d76130d73a143e76218c7c66ca0167d06f70d591f9c3482ec8e04fe895f}.
Sections~4.2.5--4.2.6 and Table~4.1, numbered pp.~105--109
(physical pp.~119--123), and Section~5.9.1, pp.~222--223.
The published edition renumbers the Cayley chapter as Section~4.3.
Author's book page:
\url{https://webspace.maths.qmul.ac.uk/r.a.wilson/TFSG.html}.

\bibitem{BT} T.~C.~Burness and A.~R.~Thomas,
\emph{The classification of extremely primitive groups},
authors' primary manuscript, Proposition~7.4(a)--(d), pp.~52--53,
\url{https://seis.bristol.ac.uk/~tb13602/docs/ep_exceptional_final.pdf}.
\end{thebibliography}
\end{document}
